\documentclass[10pt,letterpaper]{article}
\usepackage{cornell-notes}

\title{Clause 31.05.05.01: Fmtflags Manipulators}
\author{}
\date{}
\setDocTitle{Clause 31.05.05.01: Fmtflags Manipulators}
\setDocSubtitle{C++ 2024}
\setDocOwner{C++ 2024 Cornell Notes}
\setDocDate{\today}
\setCornellCollection{C++ 2024 Cornell Notes}
\setCornellUnitType{Clause}
\setCornellUnitNumber{31.05.05.01}
\setCornellUnitTitle{Fmtflags Manipulators}
\hypersetup{pdftitle={Clause 31.05.05.01: Fmtflags Manipulators},pdfsubject={Cornell notes},pdfauthor={}}

\begin{document}
\maketitle
\begin{CornellOverviewBox}{Overview}
\textbf{Classification:} Standard library - Normative\\[2pt]
\textbf{Learning objective:} Understand the rules, terminology, and semantic consequences associated with fmtflags manipulators.
\end{CornellOverviewBox}\begin{CornellSummaryBox}{Summary}
{\sffamily\bfseries\color{Accent} Summary}\par\vspace{3pt}Section 31.5.5.1 focuses on fmtflags manipulators. Apply the specified interface together with its mandates, effects, result conditions, exception behavior, and complexity constraints. When applying it, preserve the distinction between mandatory requirements, recommendations, permissions, and behavior deliberately left undefined, unspecified, or implementation-defined.
\end{CornellSummaryBox}

\end{document}
